% ECONOMICS_PROBLEM: {"id": "D44-34", "title": "Optimal auctions beyond independent private values", "jel_code": "D44", "source_id": "OP-0007"}
% FIRST_POSTED: 2026-10-09
% ATTEMPT_STATUS: unresolved
% ENTRY_KIND: research_attempt
% !TeX program = pdflatex
% Self-contained: no external input or bibliography files required.
\documentclass[11pt,a4paper]{article}
\usepackage[T1]{fontenc}
\usepackage[utf8]{inputenc}
\usepackage{lmodern}
\usepackage[margin=24mm]{geometry}
\usepackage{amsmath,amssymb,amsthm,mathtools}
\usepackage{booktabs,longtable,array,enumitem,microtype,xurl,hyperref}
\hypersetup{colorlinks=true,linkcolor=blue,citecolor=blue,urlcolor=blue,pdftitle={Checked research attempts: nine economics and stochastic-analysis problems}}
\setlength{\parindent}{0pt}
\setlength{\parskip}{3pt}
\setlength{\emergencystretch}{3em}
\widowpenalty=10000\clubpenalty=10000\displaywidowpenalty=10000
\setlist[enumerate]{leftmargin=*,itemsep=2pt,topsep=3pt}
\newcommand{\R}{\mathbb R}
\newcommand{\N}{\mathbb N}
\newcommand{\E}{\mathbb E}
\newcommand{\Pp}{\mathbb P}
\newcommand{\Q}{\mathbb Q}
\newcommand{\F}{\mathcal F}
\newcommand{\ind}{\mathbf 1}
\newcommand{\cE}{\mathcal E}
\newcommand{\Law}{\operatorname{Law}}
\newcommand{\sgn}{\operatorname{sgn}}
\newcommand{\Var}{\operatorname{Var}}
\newcommand{\dist}{\operatorname{dist}}
\newcommand{\KL}{\operatorname{KL}}
\newcommand{\TV}{\operatorname{TV}}
\newcommand{\Lip}{\operatorname{Lip}}
\DeclareMathOperator*{\esssup}{ess\,sup}
\DeclareMathOperator*{\argmax}{arg\,max}
\newtheorem{theorem}{Theorem}
\newtheorem{lemma}[theorem]{Lemma}
\newtheorem{proposition}[theorem]{Proposition}
\newtheorem{corollary}[theorem]{Corollary}
\theoremstyle{definition}
\newtheorem{example}[theorem]{Example}
\newcommand{\report}[3]{\clearpage\setcounter{theorem}{0}\renewcommand{\thetheorem}{#1.\arabic{theorem}}\renewcommand{\theHtheorem}{#1.\arabic{theorem}}\section*{#1: #2}\addcontentsline{toc}{section}{#1: #2}\textbf{Outcome:} #3.\par\textbf{Tools:} web browsing available; Python computation available. Literature checked on 9 October 2026.}
\newcommand{\stage}[2]{\subsection*{#1) #2}}
\newcommand{\unresolved}{\textbf{LABEL: UNRESOLVED.}}
\newcommand{\disproof}{\textbf{LABEL: FULL SOLUTION. SUBLABEL: COUNTEREXAMPLE/DISPROOF.}}

\begin{document}
\section*{D44-34: Robust auction design with interacting frictions}
\textbf{Outcome:} UNRESOLVED as a general research family; exact restricted results and selection obstructions proved.\par\textbf{Tools:} web browsing available; Python computation available. Literature checked on 9 October 2026.
\subsection*{1) FORMAL RESTATEMENT}
Fix finitely many bidders and objects, finite private signal sets, a declared feasible allocation set, and a collection $\Omega$ of environments. Each environment specifies the signal law, values (possibly interdependent), monetary utility, budgets, and entry or information costs. A mechanism is a finite extensive-form auction with perfect recall and specified lotteries and transfers. For every admissible $(m,\omega)$ assume a nonempty set $\mathcal E(m,\omega)$ of sequential-equilibrium assessments, with consistency defined by limits of completely mixed behavioral strategies and Bayesian beliefs. Participation, budget feasibility, cost accounting, and expectations must be declared for every admitted equilibrium, not just a preferred one.

For a feasible mechanism class $\mathfrak M_F$, put
\[
 W(m)=\inf_{\omega\in\Omega}\inf_{e\in\mathcal E(m,\omega)}
       \E_{m,\omega,e}\left[\sum_i t_i-C(x)\right],
 \qquad V^*=\sup_{m\in\mathfrak M_F}W(m).
\]
The question asks for tractable restrictions giving a characterization or approximation relative to this same feasible class. It does not specify one fixed $\Omega$ or one fixed $\mathfrak M_F$; it is a research family rather than a single universal true-or-false theorem. The minimal clarification is to supply these sets, utility and outside-option conventions, and whether participation is required at every equilibrium. We use outside utility zero and require participation in all admitted equilibria in the examples below.

The stress points are adverse equilibrium selection, ties at a limiting price, changes in entry, infeasible reports under private budgets, and the distinction between a supremum and an attained maximum. The uploaded statement already uses a supremum; a nonattainment example therefore exposes an obstacle rather than disproving that correctly phrased programme.

\subsection*{2) QUICK LITERATURE/CONTEXT CHECK}
Myerson \cite{Myerson} is a specific optimal-auction benchmark, not a universal formula for every listed friction. Brooks--Du \cite{BrooksDu} study informational robustness in a structured common-value setting. Hartline--Hoy--Taggart \cite{HHT} explicitly study worst-environment and worst-equilibrium auction performance under their own restrictions. These primary sources motivate careful separation of the uncertainty set, benchmark, and equilibrium-selection rule rather than transferring a theorem across unmatched mechanism classes.

\subsection*{3) ATTACK PLAN}
\textbf{Proof routes.} Solve small robust subfamilies exactly, formulate a finite-outcome incentive-constrained programme where valid, and prove how certified worst-equilibrium payoff enclosures yield an approximation in the same finite catalogue.

\textbf{Disproof routes.} Use a single bidder at a tie price to test attainment; put a zero-value environment in the uncertainty set to test nontriviality; construct an entry-coordination game to show that a favorable equilibrium cannot replace the infimum over equilibria.

\textbf{Landscape and tools.} Backward induction solves posted prices; participation bounds transfers; min--max inequalities give approximation bounds; lotteries linearize expected utilities; Bayesian incentive constraints encode correlation; compact polytopes yield attainment; and entry games test equilibrium selection.

\subsection*{4) WORK}
\begin{theorem}[Compact posted-price family with no robust optimizer]
There is a one-bidder, one-object instance satisfying the stated finite-game and budget requirements for which the robust supremum equals $1$ but is not attained, even though the parameter set of admissible prices is compact.
\end{theorem}
\begin{proof}
The bidder's value is the known number $1$, monetary utility is quasilinear, budget is $1$, entry cost and seller cost are zero, and $\Omega$ is a singleton. Let $\mathfrak M_F$ consist exactly of posted prices $p\in[0,1]$, followed by the bidder's accept/reject decision. Accept allocates the object and charges $p$; reject allocates nothing and charges zero. Each game is finite, with perfect recall, and all realized payments meet the budget.

For $p<1$, accepting gives utility $1-p>0$, strictly above rejection. The unique sequentially rational outcome is acceptance, and worst-equilibrium revenue is $p$. For $p=1$, both actions give zero utility. Every mixture is sequentially rational, including rejection with probability one. There is only one information set at the bidder's decision, so each of these strategies is the limit of completely mixed strategies with the same trivial beliefs and is a sequential-equilibrium assessment. Hence worst-equilibrium revenue at $p=1$ is zero. Participation holds in every case.

It follows that
\[
 W(p)=\begin{cases}p,&0\le p<1,\\0,&p=1,\end{cases}
 \qquad \sup_{p\in[0,1]}W(p)=1,
\]
but $W(p)<1$ for every admissible $p$. For every $0<\epsilon<1$, the same-class price $p=1-\epsilon$ attains $V^*-\epsilon$; it is also a $(1-\epsilon)$ approximation here since $V^*=1$. This proves both the nonattainment and its explicit approximation without changing the benchmark.
\end{proof}

\begin{proposition}[A zero-value environment forces a trivial robust optimum]
Suppose $\Omega$ contains an environment with all bidders' values identically zero, quasilinear utility, outside utility zero, and nonnegative entry costs. Suppose every admitted equilibrium satisfies ex-ante participation for every bidder, seller costs are nonnegative, and the null mechanism with zero cost is feasible. Then $V^*=0$.
\end{proposition}
\begin{proof}
In the zero-value environment, participation gives
$\E[-t_i-\text{entry cost}_i]\ge0$, hence $\E t_i\le0$ for every bidder. Thus expected seller profit is at most zero in every admitted equilibrium of every admissible mechanism in this environment. Nonemptiness of the equilibrium sets makes the robust infimum meaningful and gives $W(m)\le0$ for every $m$. The null mechanism gives profit zero in every environment and has a trivial equilibrium, so $V^*\ge0$. Combining the inequalities proves the claim. This conclusion is for the explicitly quasilinear zero-value environment; arbitrary flat monetary utility would not justify the same transfer bound.
\end{proof}

\begin{proposition}[Exact finite-outcome truthful-equilibrium programme]
Fix a finite $\Omega$, finite signal sets with full-support laws $\pi_\omega$, exogenous information, and no endogenous entry or further information acquisition. Fix a finite set $O$ of outcomes, each including a feasible allocation and a transfer vector that is budget-feasible for every possible true signal in every environment. Fix finite utility numbers $u_{i,\omega}(o,s)$ and seller profits $v_0(o)=\sum_i t_i(o)-C(x_o)$. Robust Bayesian incentive compatibility and interim participation of a direct lottery mechanism, evaluated at truthful reporting, are finite linear constraints. Maximizing worst-environment truthful-equilibrium profit over this fixed-outcome class is an attained finite linear programme whenever its feasible set is nonempty.
\end{proposition}
\begin{proof}
Let $\lambda_{r,o}$ be the probability of outcome $o$ when the report profile is $r$. Require $\lambda_{r,o}\ge0$ and $\sum_o\lambda_{r,o}=1$ for each report profile. For every environment, bidder, true signal $s_i$, and alternative report $r_i$, impose
\[
 \sum_{s_{-i}}\pi_\omega(s_{-i}\mid s_i)
 \sum_o\bigl(\lambda_{(s_i,s_{-i}),o}-\lambda_{(r_i,s_{-i}),o}\bigr)
                         u_{i,\omega}(o,s)\ge0.
\]
These are exactly the conditional expected utility comparisons under truthful reports of other bidders, including arbitrary correlation and interdependent values through the coefficients. Interim participation imposes
\[
 \sum_{s_{-i}}\pi_\omega(s_{-i}\mid s_i)
       \sum_o\lambda_{s,o}u_{i,\omega}(o,s)\ge u^{\rm out}_{i,\omega}(s_i).
\]
For example $u_{i,\omega}(o,s)=U_{i,\omega}(v_{i,\omega}(x_o,s)-t_i(o))$ is a permitted specified utility convention. Nonlinearity of $U$ does not make these constraints nonlinear in the lottery probabilities: its values are already fixed coefficients. Universal outcome budget feasibility ensures that a deviation is not surreptitiously removed or admitted according to a misreported liquidity limit.

Introduce $\rho$ and require for every $\omega$ that
\[
 \rho\le\sum_s\pi_\omega(s)\sum_o\lambda_{s,o}v_0(o).
\]
Maximize $\rho$. Every displayed constraint is linear in $\lambda,\rho$. The feasible $\lambda$ form a closed subset of a finite product of simplexes and hence a compact set. The minimum of the finitely many continuous expected-profit functions is continuous on that set. If nonempty, it attains its maximum, giving an optimal $\lambda$ and corresponding $\rho$. This proves the finite optimization and attainment assertion; input encoding and numerical accuracy must be specified before any bit-complexity claim.
\end{proof}

\textbf{Essential limitation.} Bayesian incentive compatibility certifies truthful reporting, not every equilibrium's payoff. The programme therefore does not reformulate $\inf_{e\in\mathcal E(m,\omega)}$. Transferring its value to the original objective requires an additional implementation theorem.

\begin{example}[Entry creates a good and a bad sequential equilibrium]
There are two bidders and two objects, one assigned object valued at $1$ by each bidder. Values are known, utility quasilinear, budgets are $1$, and seller costs are zero. Each bidder simultaneously chooses whether to enter. Entry costs $1/4$, paid to an outside party. If both enter, each receives its assigned object and pays the seller $1/2$. If only one enters, no object is allocated and the seller charges nothing. If neither enters, nothing happens.

Writing $N$ for no entry and $E$ for entry, the bidder payoff matrix is
\[
\begin{array}{c|cc}
 &N&E\\\hline
 N&(0,0)&(0,-1/4)\\
 E&(-1/4,0)&(1/4,1/4)
\end{array}
\]
Both $(N,N)$ and $(E,E)$ are strict Nash equilibria: an isolated entrant loses $1/4$, while with another entrant entry gains $1/4$. Encode the simultaneous choices as a finite extensive form in which the second player does not observe the first choice. Each Nash profile, together with the beliefs induced by the other player's action, is sequentially rational and is the limit of completely mixed behavioral trembles with Bayesian beliefs. Thus both are sequential equilibria. Seller profit is $0$ at $(N,N)$ and $1$ at $(E,E)$.

The option of not entering gives each bidder zero against every opponent action, so participation holds in every equilibrium, including the mixed one. All payments meet budgets. Seller profit is nonnegative for every outcome. Hence the worst-equilibrium profit is exactly zero despite the profitable equilibrium. This explicitly demonstrates the missing selection/entry step in transferring a fixed-entry truthful benchmark to the robust extensive-form objective.
\end{example}

\begin{proposition}[Same-class robust approximation from certified enclosures]
For a finite mechanism catalogue $\mathfrak M_0\subset\mathfrak M_F$ and finite $\Omega$, suppose numbers $\ell_{m,\omega}$ satisfy
\[
 \ell_{m,\omega}\le
  \inf_{e\in\mathcal E(m,\omega)}\E_{m,\omega,e}\left[\sum_i t_i-C(x)\right]
  \le\ell_{m,\omega}+\delta
\]
for every catalogue mechanism and environment. Choose
$\widehat m\in\argmax_{m\in\mathfrak M_0}\min_\omega\ell_{m,\omega}$.
Then $W(\widehat m)\ge\max_{m\in\mathfrak M_0}W(m)-\delta$.
\end{proposition}
\begin{proof}
Let $m^*$ maximize the robust value in the finite catalogue. Each environment-wise enclosure yields
$\min_\omega\ell_{m^*,\omega}\ge W(m^*)-\delta$.
The choice of $\widehat m$ makes its lower-bound minimum at least that number. Its true robust value is at least its lower-bound minimum, proving the inequality. Once the enclosures are supplied, the selection costs $O(|\mathfrak M_0||\Omega|)$ comparisons. Computing valid enclosures of all sequential-equilibrium payoffs is an explicit prerequisite, not a free oracle. This guarantee is relative to the same catalogue. Extending it to $\mathfrak M_F$ requires a separately proved approximation bound from that full class to the catalogue.
\end{proof}

\subsection*{5) VERIFICATION}
The script evaluates worst revenues at prices $0,1/2,0.9,0.99,1$, obtaining $0,1/2,0.9,0.99,0$. It enumerates the four pure entry profiles and identifies precisely $(N,N)$ and $(E,E)$ as pure equilibria. These finite checks agree with the analytical strict-deviation comparisons.

Nonattainment is consistent with the original supremum. The zero-environment argument requires quasilinearity, not merely monotonic utility. The lottery programme fixes outcomes, information and entry, uses the correct correlated conditional laws, and evaluates a selected equilibrium. The catalogue bound compares the same class on both sides; neither an adverse equilibrium nor its revenue loss is discarded.

\subsection*{6) FINAL}
\unresolved

\textbf{Strongest checked partial result.} The posted-price and zero-value instances are solved exactly. Further results give a truthful-equilibrium lottery programme, certified same-catalogue robust approximation, and an explicit entry-selection obstruction. The posted-price approximation uses the original feasible class.

\textbf{Exact first gap for the intended combined-frictions problem.} For a specified nontrivial $\Omega$ and extensive-form mechanism class, derive a tractable characterization or certified bounds for the minimum payoff over all sequential equilibria while entry, information acquisition, risk and budgets remain endogenous. The truthful incentive constraints alone do not characterize that minimum.

\textbf{Next three moves.} (1) Bound all equilibrium revenues in a finite auction catalogue, including entry. (2) Find budget-feasible implementation conditions controlling adverse equilibria. (3) Prove a catalogue-to-full-class approximation bound that accounts for tie discontinuities.

\textbf{Likely minimal obstruction.} One bidder suffices for tie-driven nonattainment; two bidders and two entry actions suffice for adverse equilibrium selection. Richer-information guarantees must first overcome these small examples.

\clearpage
\begin{thebibliography}{99}
\bibitem{Myerson}
Roger B. Myerson. \emph{Optimal Auction Design}. Mathematics of Operations Research 6(1) (1981), 58--73. DOI: 10.1287/moor.6.1.58. \url{https://pubsonline.informs.org/doi/10.1287/moor.6.1.58}.

\bibitem{BrooksDu}
Benjamin Brooks and Songzi Du. \emph{Optimal Auction Design With Common Values: An Informationally Robust Approach}. Econometrica 89(3) (2021), 1313--1360. DOI: 10.3982/ECTA16297. Author version: \url{https://benjaminbrooks.net/downloads/brooksdu_maxmin.pdf}.

\bibitem{HHT}
Jason Hartline, Darrell Hoy and Sam Taggart. \emph{Robust Analysis of Auction Equilibria}. arXiv:2310.03702v1, 5 October 2023. \url{https://arxiv.org/html/2310.03702v1}.

\end{thebibliography}
\end{document}
